% !PRINTED 328
\weberpage{328}
%<<CONT
zwischen $0$ und $\xi'$, ein anderer Theil ausserhalb dieses Intervalles liegen.

Daraus folgt nun durch Anwendung der linearen Transformation allgemein, \gesperrt{dass auf dem Kreise, der die Punkte $\alpha, x, \xi, \xi'$ trägt, wenn $\xi, \xi'$ durch (16) verknüpft sind, in jedem der beiden Theile, in die der Kreis durch $\xi$ und $\xi'$ getheilt wird, wenigstens eine der Wurzeln $\alpha$ liegt.} Der Punkt $x$ und einer der Punkte $\xi, \xi'$ sind willkürlich; der dritte Punkt ist durch (16) bestimmt.

Wir wollen insbesondere die Annahme machen, dass der Punkt $\xi'$ mit einem der beiden durch die Gleichung (12) bestimmten Punkte $X$ zusammenfalle, also etwa
\weq{17}{
X f'(x) + Y f'(y) + (Xy - Yx)\sqrt{(1-n)H} = 0
}
setzen.

Führen wir dies in (16) ein, indem wir $\xi', \eta' = X, Y$ setzen und dann den Factor $Xy - Yx$ wegheben, so ergiebt sich
\weq{18}{
(n-1)\,[\xi f'(x) + \eta f'(y)] + (\xi y - \eta x)\sqrt{(1-n)H} = 0,
}
wodurch der Punkt $\xi$ eindeutig bestimmt ist.

Gehen wir zur inhomogenen Form über, so erhalten wir [\S.\,106, (7) und (9)] aus (17) und (18)
\weq{19}{
x - X = \dfrac{n f(x)}{f'(x) + (n-1)\sqrt{\displaystyle f'(x)^2 - \frac{n}{n-1} f(x) f''(x)}},
}
\weq{20}{
x - \xi = \dfrac{n f(x)}{f'(x) + \sqrt{\displaystyle f'(x)^2 - \frac{n}{n-1} f(x) f''(x)}},
}
worin die Quadratwurzel beide Zeichen haben kann.

Die Lage der Punkte lässt sich dann so charakterisiren. Es giebt wenigstens eine Wurzel $\alpha$, so dass
\weq{21}{
\xi,\ \alpha,\ X,\ x,
}
der Grösse nach auf einander folgen, und so, dass zwischen $x$ und $X$ keine Wurzel liegt. Wir können auch $\alpha$ so auswählen, dass zwischen $\alpha$ und $X$ keine weitere Wurzel liegt. Zwischen $\alpha$ und $\xi$ können aber möglicher Weise noch andere Wurzeln liegen.
